\section{Modern Visitor Variants}
\label{sec:modern_variants}

This section explores two structural alternatives to the Classic Visitor pattern: the modern C++17 \texttt{std::variant} approach and the Acyclic Visitor pattern. The focus here is strictly on their architectural layout and operational mechanisms. A detailed comparative analysis of their advantages and trade-offs will be provided in the Discussion section.

\subsection{Modern C++ Approach: \texttt{std::variant} and \texttt{std::visit}}

\subsubsection{Overview of \texttt{std::variant}}
Introduced in the C++17 standard, \texttt{std::variant} is a type-safe union. It allows a single variable to store exactly one object from a predetermined, finite set of types. Unlike polymorphic pointers that refer to a heap-allocated base class, a \texttt{std::variant} stores its active value directly by value and maintains an internal index to track which type is currently active.

\subsubsection{Mechanism: Functors and Runtime Dispatching}
The implementation utilizing \texttt{std::variant} completely changes the dispatch mechanism by removing the abstract base class (\texttt{Shape}) and the \texttt{accept} methods. 

In this architecture, the ``Visitor'' is constructed as a Functor---a structure containing overloaded \texttt{operator()} methods for each concrete type. At runtime, the \texttt{std::visit} standard library function acts as the central dispatcher. It performs dispatch based on the active alternative stored in the \texttt{std::variant}. In common implementations, this may be realized using a dispatch table or an equivalent optimized mechanism, which automatically extracts the underlying concrete object and invokes the matching \texttt{operator()} overload within the functor.

% PLACEHOLDER FOR UML DIAGRAM
\begin{figure}[H]
    \centering
    % \includegraphics[width=0.6\textwidth]{Figures/variant_visitor_diagram.png}
    \caption{Structural representation of the \texttt{std::variant} approach, operating without base classes.}
    \label{fig:variant_diagram}
\end{figure}

\begin{lstlisting}[language=C++, caption=Modern Visitor implementation using \texttt{std::variant} and \texttt{std::visit}.]
#include <iostream>
#include <variant>
#include <vector>
#include <string>

// Independent element classes (No Base Shape required)
class Circle {
public:
    float radius;
    explicit Circle(float r) : radius(r) {}
};

class Rectangle {
public:
    float a, b;
    Rectangle(float w, float h) : a(w), b(h) {}
};

// Define the variant holding the possible types
using ShapeVariant = std::variant<Circle, Rectangle>;

// The Visitor is a Functor with overloaded call operators
struct SaveVisitor {
    std::string filePath;
    explicit SaveVisitor(const std::string& path) : filePath(path) {}

    void operator()(const Circle& c) const {
        std::cout << "Saving Circle (radius " << c.radius << ") to " << filePath << "\n";
    }
    
    void operator()(const Rectangle& r) const {
        std::cout << "Saving Rectangle to " << filePath << "\n";
    }
};

void runModernVisitor() {
    std::vector<ShapeVariant> shapes;
    shapes.push_back(Circle{5.0f});
    shapes.push_back(Rectangle{4.0f, 6.0f});

    SaveVisitor saveOp("modern_data.xml");

    for (const auto& shape : shapes) {
        // std::visit handles the runtime dispatching
        std::visit(saveOp, shape);
    }
}
\end{lstlisting}

\subsection{The Acyclic Visitor Pattern}

\subsubsection{Structural Overview}
The Acyclic Visitor pattern alters the class hierarchy by utilizing a degenerate (empty) base visitor class. Instead of a single, monolithic \texttt{Visitor} interface that declares \texttt{visit} methods for all elements, the Acyclic pattern establishes discrete, independent visitor interfaces for each concrete element (e.g., \texttt{CircleVisitor}, \texttt{RectangleVisitor}). A concrete visitor class then inherits from the degenerate base and selectively implements only the interfaces corresponding to the elements it intends to process.

% PLACEHOLDER FOR UML DIAGRAM
\begin{figure}[H]
    \centering
    % \includegraphics[width=0.7\textwidth]{Figures/acyclic_visitor_diagram.png}
    \caption{UML Class Diagram of the Acyclic Visitor Pattern.}
    \label{fig:acyclic_diagram}
\end{figure}

\subsubsection{Mechanism: Cross-Casting}
The core dispatch mechanism in the Acyclic Visitor relies on C++ Run-Time Type Information (RTTI) through cross-casting. 

When the client invokes \texttt{accept(VisitorBase\&)} on an element, the element receives the degenerate base reference. Inside the \texttt{accept} method, the concrete element attempts to execute a \texttt{dynamic\_cast} to cast the incoming base visitor reference into its own specific visitor interface pointer (e.g., \texttt{Circle} attempts to cast to \texttt{CircleVisitor*}). If the cast is successful (meaning the concrete visitor implemented that specific interface), the element calls the \texttt{visit} method. If the cast fails, the operation is safely bypassed.

\begin{lstlisting}[language=C++, caption=Acyclic Visitor Pattern leveraging cross-casting via \texttt{dynamic\_cast}.]
#include <iostream>

// 1. Degenerate Base Visitor
class VisitorBase {
public:
    virtual ~VisitorBase() = default;
};

// 2. Specific Visitor Interfaces
class Circle;
class CircleVisitor {
public:
    virtual void visit(Circle& c) = 0;
};

class Rectangle;
class RectangleVisitor {
public:
    virtual void visit(Rectangle& r) = 0;
};

// 3. Element Hierarchy
class Shape {
public:
    virtual ~Shape() = default;
    virtual void accept(VisitorBase& v) = 0;
};

class Circle : public Shape {
public:
    void accept(VisitorBase& v) override {
        // Cross-cast execution
        if (auto* cv = dynamic_cast<CircleVisitor*>(&v)) {
            cv->visit(*this);
        }
    }
};

class Rectangle : public Shape {
public:
    void accept(VisitorBase& v) override {
        // Cross-cast execution
        if (auto* rv = dynamic_cast<RectangleVisitor*>(&v)) {
            rv->visit(*this);
        }
    }
};

// 4. Concrete Visitor (Selectively implements specific interfaces)
class SaveVisitor : public VisitorBase, public CircleVisitor {
public:
    // Only implements visit for Circle. Rectangle will fail the dynamic_cast silently.
    void visit(Circle& c) override {
        std::cout << "Acyclic: Saving Circle\n";
    }
};
\end{lstlisting}